% RSE: Section 4, Results. Every number is a \metric{} macro. Order (review P5): record ->
% waterline-model accuracy and checked sectors (H2) -> audit (H1) -> grid versus observation error
% -> area versus shoreline -> ETOPO 2022 -> Great Salt Lake -> near-term emergence (H3).

\subsection{The levelled record and the emerged seabed}
\label{sec:res_record}
The levelled record spans three epochs and two directions of level change (Fig.~\ref{fig:rse_emergence}). The sea rose until \metric{levels.epoch_E1_end} and then fell, so the shoreline crossed the upper part of the shelf twice. Annual mean levels fell by \metric{rev.w.levels.annual_fall_m}~m from 1995 to 2025. The emerged band, the seabed that was water in \metric{params.highstand_label} and land in \metric{params.lowstand_label}, is widest on the northeastern coast and in Komsomolets Bay (Fig.~\ref{fig:rse_emergence}). Across sectors and levels, the observation noise floor $Q$ ranges from \metricdp{erl.text.h1.noise_p95_min_km}{1} to \metricdp{erl.text.h1.noise_p95_max_km}{1}~km.

% Figure emergence: generated by figures/rse_fig_emergence.py; placed where the text first discusses it.
% Display grid, colour mask rationale, graticule and scale bar: Supplementary section rse_si_display.tex.
\begin{figure}[!tbp]
  \centering
  \includegraphics[width=\textwidth]{../figures/out/rse_fig_emergence.pdf}
  \caption{\textbf{When the North Caspian seabed emerged.}
  (a)~Year in which each pixel of the emerged seabed became permanently dry: the first day after the last day on which the basin level stood at or above the pixel's height in the waterline elevation model (\metric{params.dem_res_m}~m pixels).
  Colours are shown only on seabed that was water in \metric{params.highstand_label} and land in \metric{params.lowstand_label}.
  The dates inherit the errors of the elevation model (Fig.~\ref{fig:rse_dem_accuracy}).
  Light blue: open sea of \metric{params.lowstand_label}. Circles: ports, piers and ferry terminals; triangles: oil and gas platforms and wells (both from OpenStreetMap); squares: cities.
  Sector names as in Fig.~\ref{fig:rse_study_area}.
  (b)~Daily basin level on the DAHITI reference (black; orange, Hydroweb shifted to DAHITI where DAHITI has gaps), with the rising limb to \metric{levels.epoch_E1_end} (E1), the early fall to \metric{params.e2_end} (E2) and the late fall after it (E3).
  The colour strip under the curve is the colour key of a, on the same time axis.
  (c)~Landsat acquisitions per calendar month in the main processing plan, the acquisitions behind every waterline.}
  \label{fig:rse_emergence}
\end{figure}


\subsection{Accuracy of the waterline model and the checked sectors}
\label{sec:res_dem}
H2 failed. Against ATL08 on the test blocks, the waterline model has an NMAD of \metric{h2.atl08.dem.nmad_m}~m, lower than the grids' \metric{erl.text.h2.atl08.product_nmad_min_m}--\metric{erl.text.h2.atl08.product_nmad_max_m}~m. Its RMSE of \metric{h2.atl08.dem.rmse_m}~m is not lower than theirs (\metric{erl.text.h2.atl08.product_rmse_min_m}--\metric{erl.text.h2.atl08.product_rmse_max_m}~m). The paired block-bootstrap intervals of both differences include zero for \metric{erl.text.h2.atl08.diff_nmad.n_ci_include_zero} of \metric{erl.text.h1.n_products} grids, against ATL06 as well as ATL08 (Table~\ref{tab:verdicts}). When the test blocks were drawn again with other seeds, H2 failed in \metric{rev.g.h2.split.n_not_supported} of \metric{rev.g.h2.split.n_seeds} draws, so failure is the typical outcome.

The failure comes from two sectors (Fig.~\ref{fig:rse_dem_accuracy}). Over all ICESat-2 blocks, the model is biased low in the Volga delta, by \metric{rev2.d.obs.atl06.volga.mean_m}~m against ATL06 (\metric{params.ci_pct}\% interval \metric{rev2.d.obs.atl06.volga.mean_lo_m} to \metric{rev2.d.obs.atl06.volga.mean_hi_m}~m, \metric{rev2.d.obs.atl06.volga.n_blocks} blocks). It is biased high in Komsomolets Bay, by \metric{rev2.d.obs.atl06.komsomolets.mean_m}~m (\metric{rev2.d.obs.atl06.komsomolets.mean_lo_m} to \metric{rev2.d.obs.atl06.komsomolets.mean_hi_m}~m, \metric{rev2.d.obs.atl06.komsomolets.n_blocks} blocks).

% Figure dem_accuracy: generated by figures/rse_fig_dem_accuracy.py; placed where the text first discusses it.
\begin{figure}[!tbp]
  \centering
  \includegraphics[width=\textwidth]{../figures/out/rse_fig_dem_accuracy.pdf}
  \caption{\textbf{Accuracy of the waterline elevation model against ICESat-2.}
  (a)~The waterline elevation model (height on the DAHITI reference, \metric{params.dem_res_m}~m, shown on a \metric{params.fig_map_res_m}~m display grid) on the emerged seabed, with the \metric{h2.atl08.n_test_blocks} test blocks of \metric{h2.block_size_km}~km that hold ICESat-2 points (vermilion squares). Checked sectors are named in bold.
  (b)~Model height against the ICESat-2 ATL08 height (orthometric height plus the datum constant of \metric{h2.atl08.datum_constant_m}~m fitted on training blocks) for the $n = \metric{h2.atl08.n_test_points}$ test-block points on bare emerged seabed; dashed line, one to one.
  Mean error \metric{h2.atl08.dem.bias_m}~m, normalised median absolute deviation (NMAD) \metric{h2.atl08.dem.nmad_m}~m and root-mean-square error \metric{h2.atl08.dem.rmse_m}~m.
  (c)~Error (model minus ICESat-2) per sector and over all test blocks.
  Black circles, the elevation model against ATL08 (filled, mean; open, median); grey squares, the same against ATL06; coloured markers, the mean error of each grid against ATL08, after the grid's vertical offset fitted on training blocks.
  Shaded, $\pm\metric{params.checked_tol_m}$~m: a sector counts as checked when all four statistics of the elevation model lie within it.
  A value beyond the axis is drawn on its edge and labelled.
  (d)~NMAD (\metric{params.nmad_scale} times the median absolute deviation from the median) of the elevation model (black) and of each grid against ATL08, on a logarithmic axis.
  Right, the number of ATL08 points and of test blocks per row; the west and interfluve sectors hold \metric{h2.atl08.n_blocks.west} test block each, so no interval can be estimated for them.}
  \label{fig:rse_dem_accuracy}
\end{figure}


Elsewhere the model is unbiased. The checked-sector criterion rested on one test block each for the Kizlyar--Dagestan and interfluve coasts. Over all blocks, the mean bias against ATL06 is \metric{rev2.d.obs.atl06.west.mean_m}~m on the Kizlyar--Dagestan coast (\metric{rev2.d.obs.atl06.west.mean_lo_m} to \metric{rev2.d.obs.atl06.west.mean_hi_m}~m, \metric{rev2.d.obs.atl06.west.n_blocks} blocks). It is \metric{rev2.d.obs.atl06.interfluve.mean_m}~m on the interfluve coast (\metric{rev2.d.obs.atl06.interfluve.mean_lo_m} to \metric{rev2.d.obs.atl06.interfluve.mean_hi_m}~m, \metric{rev2.d.obs.atl06.interfluve.n_blocks} blocks). The northeastern coast and the open shelf also lie within $\pm$\metric{params.checked_tol_m}~m, and ATL08 gives the same picture. These biases are relative to the domain-wide ICESat-2 constant, so they show that the four checked sectors agree with each other and with ICESat-2, not that the working datum is exact.

\begin{table}[!htbp]
\centering
\renewcommand{\baselinestretch}{1}\scriptsize
\renewcommand{\arraystretch}{1.1}\setlength{\tabcolsep}{3pt}
\caption{Verdicts and accuracy by grid. $c$: area-fitted offset. Sectors failed: sectors in which the grid fails the shoreline test (Eq.~\ref{eq:rule}), offset-corrected and as distributed, of seven; and offset-corrected, of the four checked sectors. Block $p$: one-sided block permutation test of larger errors in the deltaic and vegetated sectors (offset-corrected). NMAD and RMSE: against ATL08 on the test blocks, each grid with its own training-block constant; the \metric{params.ci_pct}\% intervals of the differences from the waterline model include zero for every grid.}
\label{tab:verdicts}
\begin{tabular}{lrrrrrrr}
\toprule
 & $c$ (m) & \multicolumn{2}{c}{Sectors failed (of 7)} & Checked sectors & Block $p$ & NMAD (m) & RMSE (m) \\
 & & corrected & as distributed & failed (of \metric{rev.s.checked.n_sectors}) & & & \\
\midrule
GEBCO\_2026 & \metric{h1.shape_offset.gebco2026_m} & \metric{h1.mispredicts.shape.gebco2026.n_sectors} & \metric{h1.mispredicts.total.gebco2026.n_sectors} & \metric{rev.s.item1.shape.gebco2026.validated4.n_mispredict} & \metricsf{h1.block_test.shape.gebco2026.p}{2} & \metric{h2.atl08.gebco2026.nmad_m} & \metric{h2.atl08.gebco2026.rmse_m} \\
GEBCO\_2022 & \metric{h1.shape_offset.gebco2022_m} & \metric{h1.mispredicts.shape.gebco2022.n_sectors} & \metric{h1.mispredicts.total.gebco2022.n_sectors} & \metric{rev.s.item1.shape.gebco2022.validated4.n_mispredict} & \metricsf{h1.block_test.shape.gebco2022.p}{2} & \metric{h2.atl08.gebco2022.nmad_m} & \metric{h2.atl08.gebco2022.rmse_m} \\
ETOPO~2022 & \metric{h1.shape_offset.etopo2022_m} & \metric{h1.mispredicts.shape.etopo2022.n_sectors} & \metric{h1.mispredicts.total.etopo2022.n_sectors} & \metric{rev.s.item1.shape.etopo2022.validated4.n_mispredict} & \metricsf{h1.block_test.shape.etopo2022.p}{2} & \metric{h2.atl08.etopo2022.nmad_m} & \metric{h2.atl08.etopo2022.rmse_m} \\
ETOPO1 & \metric{h1.shape_offset.etopo1_m} & \metric{h1.mispredicts.shape.etopo1.n_sectors} & \metric{h1.mispredicts.total.etopo1.n_sectors} & \metric{rev.s.item1.shape.etopo1.validated4.n_mispredict} & \metricsf{h1.block_test.shape.etopo1.p}{2} & \metric{h2.atl08.etopo1.nmad_m} & \metric{h2.atl08.etopo1.rmse_m} \\
SRTM15+ V2.7 & \metric{h1.shape_offset.srtm15p_m} & \metric{h1.mispredicts.shape.srtm15p.n_sectors} & \metric{h1.mispredicts.total.srtm15p.n_sectors} & \metric{rev.s.item1.shape.srtm15p.validated4.n_mispredict} & \metricsf{h1.block_test.shape.srtm15p.p}{2} & \metric{h2.atl08.srtm15p.nmad_m} & \metric{h2.atl08.srtm15p.rmse_m} \\
Waterline model & -- & -- & -- & -- & -- & \metric{h2.atl08.dem.nmad_m} & \metric{h2.atl08.dem.rmse_m} \\
\bottomrule
\end{tabular}
\end{table}

\subsection{Global grids misplace the observed shorelines}
\label{sec:res_h1}
The first half of H1 holds. After the area-fitted offset, \metric{erl.text.h1.n_products_majority} of the \metric{erl.text.h1.n_products} grids (\metric{erl.text.h1.majority_products}) misplace the shoreline in \metric{erl.text.h1.majority.n_sectors_min} or \metric{erl.text.h1.majority.n_sectors_max} of the seven sectors (Table~\ref{tab:verdicts}; Figs.~\ref{fig:rse_shorelines} and \ref{fig:rse_dq}). As distributed, every grid misplaces it in all seven. ETOPO~2022 fails in \metric{h1.mispredicts.shape.etopo2022.n_sectors} sector after the offset (Section~\ref{sec:res_etopo}). The grids form at most three lineages. SRTM15+ and ETOPO1 each fail all four checked sectors, and in the GEBCO lineage at least one grid fails each of them.

% Figure dq: generated by figures/rse_fig_dq.py; placed where the text first discusses it.
% Construction of the degraded-elevation-model control: Supplementary section rse_si_display.tex.
\begin{figure}[!tbp]
  \centering
  \includegraphics[width=\textwidth]{../figures/out/rse_fig_dq.pdf}
  \caption{\textbf{Grid error against observation error, by sector and grid.}
  For each grid and sector, the ratio $D/Q$ at every level at which both can be measured ($n = \metric{h1.n_evaluable.shape.gebco2026.west}$--\metric{h1.n_evaluable.shape.gebco2026.interfluve} levels per grid and sector).
  $D$ is the grid's median shoreline displacement at the level, after its area-fitted vertical offset; $Q$ is the observation noise floor at the level, the \metric{params.err_pct_hi}th percentile of the distances between the shorelines of two interleaved halves of the acquisitions.
  (a)~The four checked sectors; (b)~the three sectors whose observations are biased against ICESat-2 or have no test block.
  Dots: one level each. Boxes: \metric{params.iqr_lo_pct}th to \metric{params.iqr_hi_pct}th percentile over levels. Large markers: the median; filled where the grid fails the shoreline test in the sector ($D > Q$ at more than half of the levels), open where it passes.
  Grey, to the right of each grid: its degraded-elevation-model control (Section~\ref{sec:sep}), $D_\mathrm{control}/Q$, at the levels where both can be measured; the two GEBCO grids share one control.
  Shaded: $D$ below $Q$. Ratios of zero are drawn on the lower edge. Logarithmic axis.}
  \label{fig:rse_dq}
\end{figure}


% Figure shorelines: generated by figures/rse_fig_shorelines.py; placed where the text first discusses it.
% Choice of the level bin, waterline cells, definition of Q, graticule and scale bars: Supplementary
% section rse_si_display.tex.
\begin{figure}[!tbp]
  \centering
  \includegraphics[width=\textwidth,height=0.66\textheight,keepaspectratio]{../figures/out/rse_fig_shorelines.pdf}
  \caption{\textbf{Where the five grids put the shoreline.}
  All panels show the level bin centred at \metric{erl.hero.e3_b029.level_m}~m (DAHITI reference), observed by \metric{erl.hero.e3_b029.n_acquisitions} Landsat acquisitions in \metric{erl.hero.e3_b029.first_year}--\metric{erl.hero.e3_b029.last_year}.
  (a)~Observed sea (blue), land (beige) and cells of undetermined state (white), the observed shoreline cells (black), the zooms of c--e and the transect of b (vermilion).
  (b)~Profile along the transect on the Volga--Ural interfluve coast, which follows one ICESat-2 ATL08 beam from sea to land.
  Coloured lines: the five grids' heights after each grid's area-fitted vertical offset; black line: the waterline elevation model on the emerged seabed; grey dots: ATL08 heights; dashed line: the level.
  Ticks under the profile: where the observed shoreline (black) and each grid's predicted shoreline cross the transect; dotted vertical lines repeat the observed crossings.
  (c--e)~Zooms on the three checked coasts, one panel per grid: observed shoreline cells (black) and every predicted shoreline cell, coloured by its distance to the nearest observed shoreline cell; grey, predicted shoreline cells of unknown observed state, left out of the displacement.
  In each panel, $D$ is the grid's median displacement in the sector at this level (both directions pooled); $Q$, in the row heading, is the sector's observation noise floor at this level (Fig.~\ref{fig:rse_method}c).}
  \label{fig:rse_shorelines}
\end{figure}


The errors are largest relative to the floor on the Kizlyar--Dagestan and interfluve coasts. On the Kizlyar--Dagestan coast the four grids miss by a sector median of \metricdp{rev.s.headline.west.d_km_min}{1}--\metricdp{rev.s.headline.west.d_km_max}{1}~km against a floor of \metricdp{rev.s.headline.west.q_km}{1}~km ($D/Q$ \metric{rev.s.headline.west.dq_median_min}--\metric{rev.s.headline.west.dq_median_max}). On the interfluve coast they miss by \metricdp{rev.s.headline.interfluve.d_km_min}{1}--\metricdp{rev.s.headline.interfluve.d_km_max}{1}~km against \metricdp{rev.s.headline.interfluve.q_km}{1}~km ($D/Q$ \metric{rev.s.headline.interfluve.dq_median_min}--\metric{rev.s.headline.interfluve.dq_median_max}). On the northeastern coast the misses are \metricdp{rev.s.headline.northeast.d_km_min}{1}--\metricdp{rev.s.headline.northeast.d_km_max}{1}~km against a wider floor of \metricdp{rev.s.headline.northeast.q_km}{1}~km ($D/Q$ \metric{rev.s.headline.northeast.dq_median_min}--\metric{rev.s.headline.northeast.dq_median_max}). On the open shelf they are \metricdp{rev.s.headline.open_shelf.d_km_min}{1}--\metricdp{rev.s.headline.open_shelf.d_km_max}{1}~km against \metricdp{rev.s.headline.open_shelf.q_km}{1}~km. These verdicts held in at least \metricsf{rev.w.h1.coasts3.p_min}{2} of the level resamples.

In heights, the misses are metres. On the Kizlyar--Dagestan and interfluve coasts the four grids' displacement corresponds to \metricdp{rev2.d.vert.westinterfluve.four_Dg_min_m}{1}--\metricdp{rev2.d.vert.westinterfluve.four_Dg_max_m}{1}~m of height at the local shelf gradient, against \metricdp{rev2.d.vert.westinterfluve.Qg_min_m}{1}--\metricdp{rev2.d.vert.westinterfluve.Qg_max_m}{1}~m for the floor. In the Volga and Ural deltas and Komsomolets Bay the grids miss by \metricdp{rev.w.h1.unvalidated.d_km_min}{1}--\metricdp{rev.w.h1.unvalidated.d_km_max}{1}~km against floors of \metricdp{rev.w.h1.unvalidated.q_km_min}{1}--\metricdp{rev.w.h1.unvalidated.q_km_max}{1}~km, but there grid and observation error cannot be separated.

The second half of H1, that errors are largest in the deltaic and vegetated sectors, cannot be settled with seven sectors. The block test supports it for \metric{erl.text.h1.block.shape.significant_products} ($p$ = \metricsf{h1.block_test.shape.gebco2026.p}{2}--\metricsf{h1.block_test.shape.etopo1.p}{2}). After Holm's adjustment over the ten block tests, none of the offset-corrected tests remains significant (adjusted $p$ \metricdp{rev.w.s.item5.block_p_holm_min}{2}--\metricdp{rev.w.s.item5.block_p_holm_max}{2}), and \metric{rev.s.item5.total.block.n_significant_holm} as-distributed tests do. Blocks within a sector are not independent, and at sector level the exact test gives $p$ = \metricsf{rev.w.s.item5.sector_p_min}{2}--\metricsf{rev.w.s.item5.sector_p_max}{2}. With seven sectors that test cannot fall below \metric{rev.s.item5.sector.min_attainable_p}. We therefore regard this half as untested, not refuted.

\subsection{Grid error against observation error}
\label{sec:res_sep}
Whole-metre storage and coarse cells explain little of the error on the checked coasts. The degraded-DEM control misses the Kizlyar--Dagestan and interfluve shorelines by only \metricdp{rev.w.r.control.westinter_km_min}{1}--\metricdp{rev.w.r.control.westinter_km_max}{1}~km. The grids exceed both the floor and the control at most levels in \metric{rev.r.control.stop.n_excess_holds} of \metric{rev.r.control.stop.n_cases} grid-by-checked-sector cases. The two exceptions are GEBCO\_2022 and SRTM15+ on the open shelf. With the resolution-matched floor, each grid still fails \metric{rev.w.r.res.validated4_min} or \metric{rev.w.r.res.validated4_max} checked sectors. The rule is not oversensitive on the coasts. When the degraded-DEM control is given spatially correlated vertical errors with a \metric{rev2.oc.i5.range_km}~km range, the rule flags it in more than half of \metric{rev2.oc.i5.R.c15s.s100} realisations only once the error's standard deviation reaches \metricdp{rev2.oc.i5.c15s.west.first_sigma_m}{2}~m on the Kizlyar--Dagestan coast and \metricdp{rev2.oc.i5.c15s.interfluve.first_sigma_m}{2}~m on the interfluve coast at 15\,arc-second resolution (\metricdp{rev2.oc.i5.c1m.west.first_sigma_m}{2} and \metricdp{rev2.oc.i5.c1m.interfluve.first_sigma_m}{2}~m at 1\,arc-minute), and on the northeastern coast at 15\,arc-second resolution not within the largest error tested. These thresholds lie below the vertical equivalents of the grids' errors on the same coasts (\metricdp{rev2.d.vert.westinterfluve.four_Dg_min_m}{1}--\metricdp{rev2.d.vert.westinterfluve.four_Dg_max_m}{1}~m). On the open shelf the rule flags errors of only \metricdp{rev2.oc.i5.c15s.openshelf.first_sigma_m}{2}~m, and it flags the 1\,arc-minute control even without added error, so the open-shelf verdicts reflect grid resolution more than grid error.

The systematic level-error budget has a standard deviation of \metric{rev.w.r.budget.sigma_median_min}--\metric{rev.w.r.budget.sigma_median_max}~m. Added to the floor, it removes the open shelf for all four grids and the northeastern coast for two of them, but not the Kizlyar--Dagestan or interfluve coasts. In the Monte Carlo over the same budget, the four grids fail the Kizlyar--Dagestan, interfluve and northeastern coasts in every one of \metric{rev2.oc.i9.B} draws, and the open shelf in \metric{rev2.oc.i9.gebco2022.openshelf.p_mispredict_pct}--\metric{rev2.oc.i9.etopo1.openshelf.p_mispredict_pct}\% of them. Shifting each sector's levels by its SWOT water-surface departure changes \metric{rev.w.v.depaudit.shape_n_changed} of \metric{rev.w.v.depaudit.n_cases} verdicts, and the static channel-bank filter \metric{rev.r.static.n_changed}.

The four grids keep failing all four checked sectors under every observation and datum variant but one: with emergent vegetation counted as dry or as water, with AWEI\textsubscript{sh} or a fixed MNDWI threshold, as distributed with DAHITI's reference offset removed, without the geoid-difference field, and under both unknown-cell bounds scored against the original floor. The exception is plain. When unknown cells are measured on both sides and the floor is recomputed the same way, no grid fails any sector. That matched floor also measures the many cells whose state is unknown in the other half, so it no longer measures how precisely the shoreline is placed. It measures how much of the coast the other half could see. Unknown cells make up \metric{rev2.a.i8.share.srtm15p.interfluve.pooled_pct}--\metric{rev2.a.i8.share.etopo2022.west.pooled_pct}\% of the grids' predicted shoreline cells on these two coasts.

No better constant rescues the grids. The displacement-fitted offset lies \metric{rev.w.r.cD_below_min_m}--\metric{rev.w.r.cD_below_max_m}~m below the area-fitted one and changes \metric{rev.r.constants.n_changed_at_c_D} verdicts. Even with the best offset chosen separately in each sector, each grid fails \metric{rev.w.r.bestcase.validated4_min}--\metric{rev.w.r.bestcase.validated4_max} checked sectors. A planar offset field lowers the area error but widens the shoreline misses (Supplementary Material). Seabed change between epochs is not excluded: at the levels crossed on both limbs, the rising-limb shoreline fits the grids better in \metric{rev.s.item8.votes.k_e1_better} of \metric{rev.s.item8.votes.n} comparisons ($p$ = \metricsf{rev.s.item8.votes.p}{2}).

Three checks were not run, because each needs new image processing: a split of the scenes by wind regime, Sentinel-2 in the lowest level bins, and the levelling of Landsat scenes from before 1992 with gauge records.

\subsection{Matching the flooded area does not place the shoreline}
\label{sec:res_area}
As distributed, every grid floods too much sea: its median area bias is \metric{erl.text.h1.area_bias_total_min_pct}--\metric{erl.text.h1.area_bias_total_max_pct}\%. The area-fitted offsets of \metric{erl.text.h1.shape_offset_min_m} to \metric{erl.text.h1.shape_offset_max_m}~m remove most of it, leaving median biases within \metric{erl.text.h1.area_bias_shape_absmax_pct}\% (Fig.~\ref{fig:rse_area}). The negative offsets mean that the grids' seabed sits too deep relative to the levelled waterlines.

% Figure area: generated by figures/rse_fig_area.py; placed where the text first discusses it.
\begin{figure}[!tbp]
  \centering
  \includegraphics[width=\textwidth]{../figures/out/rse_fig_area.pdf}
  \caption{\textbf{Agreement in flooded area is not agreement in shoreline position.}
  (a)~Area bias of each grid's predicted sea in the study domain, (predicted $-$ observed) as a percentage of the observed sea, at every level and period of the record (\metric{rev.s.item4.area_bias.shape.gebco2026.domain.n_units} units per grid), after each grid's vertical offset fitted to flooded area (\metric{rev.w.s.oos.c_all_m_min} to \metric{rev.w.s.oos.c_all_m_max}~m).
  One panel per period: rising limb (E1), early fall (E2) and late fall (E3). Errors on cells of known state outside the observed area are counted.
  Values beyond the axis are drawn on its edge (open triangles) and listed in the panel.
  (b)~Out-of-sample test of the offset: area bias in the late fall, in the calibration and test blocks, with the offset fitted on all periods' training blocks (open; \metric{rev.w.s.oos.c_all_m_min} to \metric{rev.w.s.oos.c_all_m_max}~m) and with the offset fitted on the rising limb and early fall only (filled; \metric{rev.w.s.oos.c_e12_m_min} to \metric{rev.w.s.oos.c_e12_m_max}~m).
  Markers: the median over the $n = \metric{rev.s.item4.oos.gebco2026.calibtest.c_e12.n_units}$ late-fall levels; bars: \metric{params.iqr_lo_pct}th to \metric{params.iqr_hi_pct}th percentile.
  (c)~Three ways of fitting a grid's vertical datum to flooded area on the training blocks: one offset for the domain (the main analysis), one offset per sector, and a planar offset field (an offset varying linearly with easting and northing).
  Left: the summed absolute area error over the training-block units, the quantity each fit minimises; the planar field lowers it by \metric{rev.w.v.planar.reduction_pct_min}--\metric{rev.w.v.planar.reduction_pct_max}\% against one offset.
  Right: the median over the seven sectors of each sector's median shoreline displacement.
  The planar field was audited for GEBCO\_2026 and ETOPO1 only, whose displacement rises from \metric{rev.w.v.planar.gebco2026.before_km} to \metric{rev.w.v.planar.gebco2026.after_km}~km and from \metric{rev.w.v.planar.etopo1.before_km} to \metric{rev.w.v.planar.etopo1.after_km}~km.}
  \label{fig:rse_area}
\end{figure}


One offset cannot fix the shoreline, because the offset needed depends on the level. Fitted in half-metre bands, it runs from \metric{rev2.d.const.across.band_m29p0_m28p5.max_m} to \metric{rev2.d.const.across.band_m29p0_m28p5.min_m}~m in the lowest band to \metric{rev2.d.const.across.band_m26p0_m25p5.max_m} to \metric{rev2.d.const.across.band_m26p0_m25p5.min_m}~m in the highest. Its slope is \metric{rev2.d.const.gebco2026.band_slope}~m per metre of level for both GEBCO grids, \metric{rev2.d.const.srtm15p.band_slope} for SRTM15+, \metric{rev2.d.const.etopo1.band_slope} for ETOPO1 and \metric{rev2.d.const.etopo2022.band_slope} for ETOPO~2022. Where two epochs share levels, their offsets differ by little, so the apparent change between epochs is a change with level. The grids therefore misrepresent the slope of the shelf. Out of sample, an offset fitted on the rising limb and early fall leaves median area biases of \metric{rev.s.item4.oos.calibtest.c_e12.median_min_pct} to \metric{rev.s.item4.oos.calibtest.c_e12.median_max_pct}\% in the late fall, against \metric{rev.w.s.oos.c_all_pct_min} to \metric{rev.w.s.oos.c_all_pct_max}\% with the all-epoch offset on the same cells.

The grids err together. Within the checked sectors, the median correlation between grids of the signed displacement at each level is \metric{rev2.d.corr.northeast.pearson_median}--\metric{rev2.d.corr.interfluve.pearson_median} (Pearson), and \metric{rev2.d.corr.northeast.pearson_cross_lineage_median}--\metric{rev2.d.corr.interfluve.pearson_cross_lineage_median} between grids of different lineages. They share data. GEBCO\_2026 equals GEBCO\_2022 in \metric{rev2.d.same.g26_g22.domain_pct}\% of the domain's cells and \metric{rev2.d.same.g26_g22.band_pct}\% of the emerged band. ETOPO~2022 equals GEBCO\_2022 in \metric{rev2.d.same.etopo_g22.lake_domain_pct}\% of its lake-depth cells. In GEBCO\_2026's type-identifier grid, \metric{rev2.d.tid.band.predicted_pct}\% of the emerged band is coded as depth predicted from satellite-derived gravity and \metric{rev2.d.tid.band.pregenerated_pct}\% as a pre-generated grid of mixed sources \citep{gebco2026}. The remaining \metric{rev2.d.tid.band.tid44_pct}\% is coded as gridded depths from multiple measured and derived sources, interpolated with gravity guidance. No cell of the band is coded as a direct measurement.

\subsection{ETOPO 2022 fits where it holds land heights measured after emergence}
\label{sec:res_etopo}
ETOPO~2022 passes the checked sectors for a reason that does not involve lake depths. Inside the North Caspian it holds GEBCO\_2022's depths capped at \metric{rev.w.g.etopo.cap_m}~m within the February 2000 lake outline, and FABDEM land heights outside it (Section~\ref{sec:grids}). Its per-level $\log(D/Q)$ changes at one level, \metric{rev2.d.cp.etopo2022.edge_m}~m, found by a single changepoint fit ($R^2$ = \metric{rev2.d.cp.etopo2022.r2}). Above that level its median $D/Q$ is \metric{rev2.d.cp.etopo2022.dq_median_above}, and \metric{rev2.d.cp.etopo2022.share_d_gt_q_above_pct}\% of sector-levels exceed the floor. Below it the median $D/Q$ is \metric{rev2.d.cp.etopo2022.dq_median_below}, and \metric{rev2.d.cp.etopo2022.share_d_gt_q_below_pct}\% exceed it. Above the break, most of its predicted shoreline lies on the edge of the 2000 outline or on FABDEM. Below it, the shoreline follows GEBCO\_2022's depths.

Its pass is fragile. As distributed it fails all seven sectors, and with DAHITI's reference offset removed it still fails \metric{rev2.a.i7.total_cref.etopo2022.n_validated4} checked sectors. With emergent vegetation counted as dry it fails \metric{rev2.a.i11.vegasdry.etopo2022.n_validated4}. In the Monte Carlo over the level-error budget it fails the interfluve coast in \metric{rev2.oc.i9.etopo2022.interfluve.p_mispredict_pct}\% of draws. A land surface measured after the seabed emerged can pass this test, but that is not evidence that the grid holds lake depths.

\subsection{Great Salt Lake: the scale of a passing surface}
\label{sec:res_gsl}
None of the grids holds lake-bed data at Great Salt Lake. GEBCO\_2026 codes \metric{rev2.d.tid.gsl.land_pct}\% of the band in both arms as land, and the SRTM-derived grids hold one flat height on \metric{rev.w.g.gsl.flat_share_pct_min}--\metric{rev.w.g.gsl.flat_share_pct_max}\% of it. Every grid fails both arms, so the pre-specified verdict is ``\metric{generality.great_salt_lake.verdict}'' (Fig.~\ref{fig:rse_gsl}). Because the grids hold no lake-bed data there, we read this as a test of the method, not as evidence that the Caspian errors generalise.

% Figure gsl: generated by figures/rse_fig_gsl.py; placed where the text first discusses it.
\begin{figure}[!tbp]
  \centering
  \includegraphics[width=\textwidth]{../figures/out/rse_fig_gsl.pdf}
  \caption{\textbf{The same audit at Great Salt Lake, Utah, with a lake-bed control from echo soundings.}
  (a, b)~For the north and south arms, the shoreline displacement $D$ at each level bin (levels on NAVD~88), after each surface's vertical offset fitted to flooded area, against the observation noise floor $Q$ (grey area; the \metric{params.err_pct_hi}th percentile of the distances between the shorelines of two interleaved halves of the acquisitions).
  Coloured lines: the five global grids; open markers on the top edge, levels at which a grid predicts no shoreline in the arm.
  Grey dashed line: the 3DEP control, mostly 2016 lidar.
  Black line: the lake-bed control, built from USGS single-beam echo soundings below the sounded top and from digitised map contours and the National Elevation Dataset above it (offset \metric{rev2.g.lakebed.all.north.shape_offset_m}~m north and \metric{rev2.g.lakebed.all.south.shape_offset_m}~m south); filled where $D > Q$ at the level.
  Shaded columns: the sounded levels, bins \metric{rev2.g.lakebed.sounded.north.shape.bins} in the north arm (\metric{rev2.g.lakebed.coverage.north.sounded.level_low_navd88_m}--\metric{rev2.g.lakebed.coverage.north.sounded.level_high_navd88_m}~m) and \metric{rev2.g.lakebed.sounded.south.shape.bins} in the south arm (\metric{rev2.g.lakebed.coverage.south.sounded.level_low_navd88_m}--\metric{rev2.g.lakebed.coverage.south.sounded.level_high_navd88_m}~m).
  With the offset fitted over all its levels, the lake-bed control exceeds $Q$ at \metric{rev2.g.lakebed.all.north.class.sounded.n_exceeding} of \metric{rev2.g.lakebed.all.north.class.sounded.n_evaluable} sounded levels in the north arm and \metric{rev2.g.lakebed.all.south.class.sounded.n_exceeding} of \metric{rev2.g.lakebed.all.south.class.sounded.n_evaluable} in the south arm.
  (c)~For each arm and surface, the median of $D$ over the evaluable levels (marker) against the median $Q$ over the same levels (grey bar ending in a tick); filled where the surface fails the shoreline test in the arm ($D > Q$ at more than half of the levels).
  The lake-bed control is shown over all its levels (\metric{rev2.g.lakebed.all.north.shape.n_evaluable} north, \metric{rev2.g.lakebed.all.south.shape.n_evaluable} south) and over the sounded levels alone, with its offset refitted on them; on the sounded levels it fails the test in both arms (\metric{rev2.g.lakebed.sounded.north.shape.n_exceeding} of \metric{rev2.g.lakebed.sounded.north.shape.n_evaluable} and \metric{rev2.g.lakebed.sounded.south.shape.n_exceeding} of \metric{rev2.g.lakebed.sounded.south.shape.n_evaluable} levels), so the control itself misplaces the sounded shorelines.
  Logarithmic axes.
  (d)~Heights of each grid over each arm's emerged band (share of the band's cells per half-metre bin of height); dotted line, the gauge level of February 2000 (\metric{rev.g.gsl.north.gauge_feb2000_navd88_m}~m north, \metric{rev.g.gsl.south.gauge_feb2000_navd88_m}~m south), when SRTM flew.}
  \label{fig:rse_gsl}
\end{figure}


The test resolves differences of a few hundred metres. The 3DEP lidar control misses by a median of \metricsf{generality.great_salt_lake.control.north.median_displacement_m}{2}~m in the north arm and \metricsf{generality.great_salt_lake.control.south.median_displacement_m}{2}~m in the south arm, and passes. The sounding-based lake-bed model \citep{tarboton_gsl_bathymetry} passes over all levels. At the levels where it rests on the sonar surveys \citep{baskin2005,baskin2006}, it misses by \metricsf{rev2.g.lakebed.sounded.south.shape.median_displacement_m}{2}--\metricsf{rev2.g.lakebed.sounded.north.shape.median_displacement_m}{2}~m and fails both arms. At those levels the grids miss by \metricsf{rev2.g.lakebed.grids_at_sounded.gebco2026.south.median_displacement_m}{2}~m (GEBCO\_2026, south arm) to \metricsf{rev2.g.lakebed.grids_at_sounded.gebco2022.north.median_displacement_m}{2}~m (GEBCO\_2022, north arm). Over all levels the grids' displacement is \metric{rev2.g.lakebed.ratio.etopo2022.north.median}--\metric{rev2.g.lakebed.ratio.srtm15p.south.median} times the lake-bed model's. The test flags even a survey made years before the water reached these levels, at a few hundred metres. The grids miss by kilometres.

\subsection{Near-term emergence}
\label{sec:res_nearterm}
Re-referenced to the observed 2026 shoreline, the grids place \metricsf{rev2.n.reref.area.main.E.1.domain.min_km2}{2}--\metricsf{rev2.n.reref.area.main.E.1.domain.max_km2}{3}~\kmsq{} of the observed sea in the first metre of further fall (Fig.~\ref{fig:rse_nearterm}). The GEBCO grids give the least and SRTM15+ the most. The re-referencing constants $c_{26}$ are \metric{rev2.n.reref.area.c26.min_m}--\metric{rev2.n.reref.area.c26.max_m}~m, so the hindcast offset does not carry over to the 2026 shoreline. Each grid still calls \metricsf{rev2.n.reref.area.main.M.min_km2}{2}--\metricsf{rev2.n.reref.area.main.M.max_km2}{2}~\kmsq{} of the observed 2026 sea dry (\metricsf{rev2.n.reref.area.main.M.min_pct_of_S_obs}{2}--\metricsf{rev2.n.reref.area.main.M.max_pct_of_S_obs}{3}\%). Moving $L_{2026}$ by $\pm$\metric{rev.w.n.sens.jitter_m}~m, or fitting $c_{26}$ to the three lowest levels, widens the first-metre range to \metricsf{rev2.n.reref.allsens.E.1.domain.min_km2}{2}--\metricsf{rev2.n.reref.allsens.E.1.domain.max_km2}{3}~\kmsq{}. The half-metre estimate is not robust to the choice of $c_{26}$. Declines beyond \metric{rev.n.env.max_pair_decline_m}~m, the largest in the hindcast, are uncalibrated; past errors of emerged area within that range are in the Supplementary Material.

% Figure nearterm: generated by figures/rse_fig_nearterm.py; placed where the text first discusses it.
% Display block, shading construction, clipping and scale bar: Supplementary section rse_si_display.tex.
% The first-metre range of the five grids is given in the Results text (sec:res_nearterm).
\begin{figure}[!tbp]
  \centering
  \includegraphics[width=\textwidth]{../figures/out/rse_fig_nearterm.pdf}
  \caption{\textbf{The next metres, measured from the observed 2026 shoreline.}
  Each grid is offset so that its sea at the 2026 level (\metric{h3.level_2026_m}~m, DAHITI reference) has the area of the observed 2026 sea (\metricsf{rev2.n.reref.S_obs_km2}{3}~km\textsuperscript{2}); the observed sea that a grid still places on dry ground (the misfit) is \metricsf{rev2.n.reref.area.main.M.min_km2}{2}--\metricsf{rev2.n.reref.area.main.M.max_km2}{2}~km\textsuperscript{2} (\metricdp{rev2.n.reref.area.main.M.min_pct_of_S_obs}{1}--\metricdp{rev2.n.reref.area.main.M.max_pct_of_S_obs}{1}\%).
  (a)~Number of grids that place each cell of the observed 2026 sea in the first metre of further fall (uncoloured, none).
  At least one grid places \metricsf{rev2.n.agree60.observed_reref.1.any_km2}{3}~km\textsuperscript{2} there and all five agree on \metricsf{rev2.n.agree60.observed_reref.1.all_km2}{2}~km\textsuperscript{2} (\metricdp{rev2.n.agree60.observed_reref.1.all_pct_of_any}{1}\%).
  Light blue: the open sea of \metric{params.lowstand_label} from Landsat.
  (b)~Emerging seabed $E(D)$, the area that emerges between the 2026 shoreline and a further fall $D$, for each grid: solid up to \metric{rev.n.env.max_pair_decline_m}~m, the largest fall in the hindcast, dotted beyond (not calibrated).
  Shading: the range of $E(D)$ when the 2026 level is raised or lowered by \metric{rev.w.n.sens.jitter_m}~m.
  (c)~Fall below the 2026 level from 2026 to 2070 for the seven primary published pathways (each pathway's own change after 2026); grey band, the range pathways (inter-model ranges and quartiles); dashed line, the calibrated limit of b.}
  \label{fig:rse_nearterm}
\end{figure}


As distributed, the grids put \metric{rev.w.n.dry.unc1_pct_min}--\metric{rev.w.n.dry.unc1_pct_max}\% of their first metre on seabed that was already dry in \metric{params.lowstand_label}. That is a direct consequence of the offsets of Section~\ref{sec:res_area}.

The grids agree on where the first metre emerges more than independent errors would allow: all five place a fraction \metricsf{rev2.n.sim.agreement_observed_reref_1m.observed_480m}{2} of the re-referenced band in it, against \metricsf{rev2.n.sim.agreement_observed_reref_1m.p50.min}{2}--\metricsf{rev2.n.sim.agreement_observed_reref_1m.p50.max}{2} for five perturbed copies of one grid. They share sources, so this agreement is not evidence of accuracy.

The pathway decides how much emerges by mid-century. On the re-referenced baseline, the pathway carries a share of \metric{rev2.n.var.observed_reref.products.domain.2035.area.pathway} of the variance of emerged area at 2035 and the grid \metric{rev2.n.var.observed_reref.products.domain.2035.area.grid}. At 2050 the shares are \metric{rev2.n.var.observed_reref.products.domain.2050.area.pathway} and \metric{rev2.n.var.observed_reref.products.domain.2050.area.grid}. Of the \metric{erl.text.h3.paths.n_primary} primary pathways, \metric{erl.text.h3.paths.1m.n_reached} reach the first metre between \metric{erl.text.h3.paths.1m.first_year} and \metric{erl.text.h3.paths.1m.last_year}; the 2016--2025 trend reaches it by \metric{rev.w.n.trend.2016_2025.year_1m}.

Exposure follows where each grid puts the band. Without the WDPA polygon that is drawn far larger than its reserve, protected areas hold \metric{rev2.n.exp.wdpa_protected_areas_noksrz.observed_reref.1.min_km2}--\metricsf{rev2.n.exp.wdpa_protected_areas_noksrz.observed_reref.1.max_km2}{2}~\kmsq{} of the first metre (a share of \metric{rev2.n.exp.wdpa_protected_areas_noksrz.observed_reref.1.min_share}--\metric{rev2.n.exp.wdpa_protected_areas_noksrz.observed_reref.1.max_share}); with it, \metric{rev2.n.exp.wdpa_protected_areas.observed_reref.1.min_share}--\metric{rev2.n.exp.wdpa_protected_areas.observed_reref.1.max_share}. The first metre exposes \metricdp{rev2.n.exp.osm_channels.observed_reref.1.min}{1}--\metricdp{rev2.n.exp.osm_channels.observed_reref.1.max}{1}~km of mapped channels and cuts \metricdp{rev2.n.exp.osm_ports.observed_reref.1.min}{0}--\metricdp{rev2.n.exp.osm_ports.observed_reref.1.max}{0} ports or landings off from the sea. The settlement count, \metricdp{rev2.n.exp.osm_settlements.observed_reref.1.min}{0}--\metricdp{rev2.n.exp.osm_settlements.observed_reref.1.max}{0}, falls to \metricdp{rev2.n.exp.osm_settlements.observed_reref_misfit_sea.1.max}{0} for every grid if misfit cells are treated as sea, so it is not robust.

H3, as specified, is ``\metric{h3.verdict.etopo2022}'' for ETOPO~2022 and ``\metric{h3.verdict.gebco2026}'' for the other four grids. The verdict is an artefact. With grids stored in whole metres, the Theil--Sen error model becomes very steep, and it applies shifts of up to \metric{h3.atl24.error_model.max_applied_shift_m}~m, outside the fitted range in \metric{h3.atl24.error_model.n_outside_fit_range} of \metric{h3.atl24.error_model.n_models} grid--sector models. Fitted without the lowest epoch, the correction does worse than no correction there for \metric{erl.text.h3.registered.e3_n_corrected_worse} of \metric{erl.text.h1.n_products} grids. The planned independent check with ATL24 bathymetry could not be made (Section~\ref{sec:heights}).
